\chapter{Elementary Reductions of the Riemann Hypothesis}
\label{v4-arith:w5-direct:closure-path}

The arithmetic Koszul reductions isolate the zero-placement problem
inside the Weil distribution.  The classical equivalences are recalled
with their hypotheses, and the restricted Paley--Wiener class is kept
separate from the full Weil cone.

\section{The Riemann Hypothesis in Classical Form}
\label{v4-arith:w5-direct:RH-classical}

\begin{definition}[the zeta function]
\label{v4-arith:w5-direct:def:zeta}
For \(\mathrm{Re}(s)>1\), the Riemann zeta function is the absolutely
convergent sum
\begin{equation}
\label{v4-arith:w5-direct:eq:zeta-dirichlet}
\zeta(s)\;:=\;\sum_{n=1}^{\infty}\frac{1}{n^{s}}.
\end{equation}
Its Euler product, equivalent on \(\mathrm{Re}(s)>1\), is
\begin{equation}
\label{v4-arith:w5-direct:eq:zeta-euler}
\zeta(s)\;=\;\prod_{p}\bigl(1-p^{-s}\bigr)^{-1},
\end{equation}
the product taken over all prime numbers.
\end{definition}

\begin{definition}[completed zeta and xi]
\label{v4-arith:w5-direct:def:xi}
Set
\begin{align}
\label{v4-arith:w5-direct:eq:Lambda}
\Lambda_{\zeta}(s)
&\;:=\;\pi^{-s/2}\,\Gamma(s/2)\,\zeta(s),
\\[2pt]
\label{v4-arith:w5-direct:eq:xi}
\xi_{\mathbb{R}}(s)
&\;:=\;\tfrac{1}{2}\,s(s-1)\,\Lambda_{\zeta}(s).
\end{align}
The function \(\Lambda_{\zeta}\) is meromorphic with simple poles at
\(s=0\) and \(s=1\); \(\xi_{\mathbb{R}}\) is entire. Both satisfy the
functional equation
\begin{equation}
\label{v4-arith:w5-direct:eq:functional-equation}
\Lambda_{\zeta}(s)\;=\;\Lambda_{\zeta}(1-s),
\qquad
\xi_{\mathbb{R}}(s)\;=\;\xi_{\mathbb{R}}(1-s).
\end{equation}
\end{definition}

\begin{remark}[functional equation via Poisson summation]
\label{v4-arith:w5-direct:rem:Poisson}
\eqref{v4-arith:w5-direct:eq:functional-equation} admits a
two-line derivation. Let \(\theta(t):=\sum_{n\in\mathbb{Z}}e^{-\pi n^{2}t}\)
for \(t>0\). Poisson summation gives
\(\theta(1/t)=t^{1/2}\theta(t)\). Integrating
\((\theta(t)-1)/2\) against \(t^{s/2-1}\,dt/2\) on \([0,\infty)\) produces
\(\Lambda_{\zeta}(s)\); the change of variable \(t\mapsto 1/t\) converts
the integral into its \((1-s)\)-partner and supplies the residual
polar term \(\tfrac{1}{2}(s(s-1))^{-1}\) absorbed into \(\xi_{\mathbb{R}}\).
Sources: B.~Riemann, ``\"{U}ber die Anzahl der Primzahlen unter einer gegebenen Gr\"{o}\ss e'', Monatsber.~Berlin.~Akad.~(1859); E.~C.~Titchmarsh, ``The Theory of the Riemann Zeta-Function'', Oxford 1951, Ch.~2.
\end{remark}

\begin{conjecture}[Riemann Hypothesis]
\label{v4-arith:w5-direct:conj:RH}
All non-trivial zeros of \(\zeta\) have real part \(1/2\). Equivalently,
the zeros of \(\xi_{\mathbb{R}}\) have real part \(1/2\).
\end{conjecture}

\section{Four Classical Equivalents}
\label{v4-arith:w5-direct:four-equivalents}

Four classical restatements of
\Cref{v4-arith:w5-direct:conj:RH} are used below.  They do not
identify the same object; they identify the same zero-placement
condition through different test categories.

\begin{theorem}[Weil positivity]
\label{v4-arith:w5-direct:thm:weil-positivity}
\ClaimStatusProvedElsewhere. Let
\(\mathrm{PW}(\mathbb{R})\) denote the additive Paley--Wiener test class
on the logarithmic line.  For \(f\in\mathrm{PW}(\mathbb{R})\), put
\[
\Phi_f(s):=\int_{\mathbb{R}}f(u)e^{(s-1/2)u}\,du
          =\int_{\mathbb{R}_{>0}}\phi(x)x^{s-1/2}\,d^{\times}x,
\qquad \phi(e^{u})=f(u).
\]
Thus \(\Phi_f\) is the Mellin transform on \(\mathbb{R}_{>0}^{\times}\)
written in logarithmic coordinates.  Define the Weil
distribution
\begin{equation}
\label{v4-arith:w5-direct:eq:weil-distribution}
W(f)\;:=\;\Phi_f(0)+\Phi_f(1)
-\sum_{\rho}\Phi_f(\rho)
-\int_{-\infty}^{\infty}\Phi_f(1/2+it)
\,\frac{\psi(1/4+it/2)+\log\pi}{2\pi}\,dt
-2\sum_{p,n}\frac{\log p}{p^{n/2}}f(n\log p),
\end{equation}
where the inner sum is over non-trivial zeros \(\rho\) of \(\zeta\) and
\(\psi=\Gamma'/\Gamma\). Then RH holds if and only if
\(W(f\ast\widetilde f)\geq 0\) for all \(f\in\mathrm{PW}(\mathbb{R})\),
where \(\widetilde f(x):=\overline{f(-x)}\). Source: A. Weil, ``Sur les
``formules explicites'' de la th\'eorie des nombres premiers'',
Meddelanden Lund 1952 (Oeuvres Scientifiques II).
\end{theorem}

\begin{theorem}[Li positivity]
\label{v4-arith:w5-direct:thm:li-positivity}
\ClaimStatusProvedElsewhere. Define
\begin{equation}
\label{v4-arith:w5-direct:eq:li-coef}
\lambda_{n}\;:=\;\sum_{\rho}\left(1-\Bigl(1-\frac{1}{\rho}\Bigr)^{n}\right)
\;=\;\frac{1}{(n-1)!}\frac{d^{n}}{ds^{n}}
\bigl[s^{n-1}\log\xi_{\mathbb{R}}(s)\bigr]\Big|_{s=1}.
\end{equation}
Then RH holds if and only if \(\lambda_{n}\geq 0\) for all \(n\geq 1\).
Source: X.-J. Li, ``The positivity of a sequence of numbers and the
Riemann hypothesis'', J.~Number Theory 65 (1997), 325--333.
\end{theorem}

\begin{theorem}[Hilbert--P\'olya spectral form]
\label{v4-arith:w5-direct:thm:HP-form}
\ClaimStatusProvedElsewhere. If there exists a separable Hilbert space
\(\mathcal{H}\) and a self-adjoint (densely defined) operator \(H\) on
\(\mathcal{H}\) such that
\(\Theta=\tfrac12\mathrm{id}+iH\) has regularised determinant
\(\xi_{\mathbb R}\), then RH holds. Equivalently, the spectral
parameters \(z_{\rho}=-i(\rho-\tfrac12)\) are all real. Sources: D.~Hilbert,
G.~P\'olya, folklore; A.~Connes, ``Trace formula in noncommutative
geometry and the zeros of the Riemann zeta function'', Selecta Math.~5
(1999), arXiv:math/9811068.
\end{theorem}

\begin{theorem}[de Branges form]
\label{v4-arith:w5-direct:thm:dB-form}
\ClaimStatusProvedElsewhere. Let
\(\Xi(z):=\xi_{\mathbb R}(\tfrac12+iz)\).  A de Branges datum for
\(\Xi\) is an entire function \(E_{\xi}\) of finite exponential type
such that \(E_{\xi}\) is Hermite--Biehler and the associated divisor
map recovers the zero divisor of \(\Xi\).  The existence of such a
datum in the zeta case is equivalent to RH.  The general
Hermite--Biehler theory is de Branges' 1968 theorem; the zeta-specific
reformulation is due to de Branges 1994 and Lagarias 1999.
\end{theorem}

\begin{remark}[equivalence of the four]
\label{v4-arith:w5-direct:rem:four-equivalents}
\Cref{v4-arith:w5-direct:thm:weil-positivity},
\Cref{v4-arith:w5-direct:thm:li-positivity},
\Cref{v4-arith:w5-direct:thm:HP-form}, and
\Cref{v4-arith:w5-direct:thm:dB-form}
are all equivalent to RH. They are not equivalent to each other in an
obvious elementary fashion: Weil positivity trades over test
functions; Li positivity trades over coefficients; Hilbert--P\'olya
trades over operators; de~Branges trades over entire functions. The
arithmetic reductions below compare these criteria after the relevant
trace, determinant, and domain data have been fixed.
\end{remark}

\section{The Arithmetic Reduction}
\label{v4-arith:w5-direct:programme-elementary}

The arithmetic Koszul construction supplies five comparison objects.
Each object is defined on its own domain, and no implication from a
restricted domain to the full Weil cone is used without an extension
theorem.

\subsection{Step 1: The arithmetic Heisenberg Koszul pair}
\label{v4-arith:w5-direct:step-1}

\begin{definition}[arithmetic Heisenberg algebra]
\label{v4-arith:w5-direct:def:heisenberg}
Let \(\mathsf{H}^{\mathrm{Ar}}\) be the completed tensor product
\begin{equation}
\label{v4-arith:w5-direct:eq:heis}
\mathsf{H}^{\mathrm{Ar}}\;:=\;
\mathsf{H}_{\infty}\;\widehat\otimes\;\widehat\bigotimes_{p}\,\mathsf{H}_{p},
\end{equation}
where \(\mathsf{H}_{\infty}\) is the standard chiral Heisenberg vertex
algebra on \(\mathbb{C}_{\infty}\) and each \(\mathsf{H}_{p}\) is the
completed Iwasawa group algebra of \(G_{p}=\mathrm{Gal}(\mathbb{Q}_{p}^{\mathrm{ab}}/\mathbb{Q}_{p})\)
with generator \(\alpha_{p}\) corresponding to Frobenius lift. The
completion is the Tate restricted tensor product of
\Cref{v4-arith:tate:setup:data}.
\end{definition}

\begin{definition}[arithmetic Koszul functor]
\label{v4-arith:w5-direct:def:koszul}
Let \(\kappa^{\mathrm{Ar}}_{\mathsf{H}}\) be the arithmetic bar complex of
\(\mathsf{H}^{\mathrm{Ar}}\) and \((\kappa^{!})^{\mathrm{Ar}}_{\mathsf{H}}\)
its arithmetic cobar. These are functors from chain complexes of
\(\mathsf{H}^{\mathrm{Ar}}\)-modules to chain complexes of
\((\mathsf{H}^{\mathrm{Ar}})^{!}\)-comodules and back; their composition
is the identity on modules up to quasi-isomorphism whenever the
arithmetic Positselski compatibility holds
(\Cref{v4-arith:thm:arithmetic-positselski}).
\end{definition}

\begin{theorem}[P1: functional equation from Koszul self-duality]
\label{v4-arith:w5-direct:thm:P1-elementary}
\ClaimStatusProvedHere. The identity
\(\kappa^{\mathrm{Ar}}_{\mathsf{H}}+(\kappa^{!})^{\mathrm{Ar}}_{\mathsf{H}}=0\),
interpreted as the cancellation of the bar and cobar differentials on
the primary diagonal subspace, implies
\(\xi_{\mathbb{R}}(s)=\xi_{\mathbb{R}}(1-s)\).
Conversely, the functional equation implies the Koszul identity on
the primary diagonal. This is
\Cref{v4-arith:thm:G-row-Riemann-functional-equation}.
\end{theorem}

\begin{remark}[elementary content of P1]
\label{v4-arith:w5-direct:rem:P1-elementary}
The primary diagonal of \(\mathsf{H}^{\mathrm{Ar}}\) carries a character
\(\chi_{s}\) defined by \(\alpha_{p}\mapsto p^{-s}\) and
\(\alpha_{\infty}\mapsto e^{i\pi s/2}\Gamma(s/2)\). The bar complex
pairs \(\chi_{s}\) with its cobar dual \(\chi_{s}^{!}\); the dual is
\(\chi_{1-s}\) by Poisson summation and Tate's local functional
equation at each place. The Koszul cancellation statement
\(\kappa+\kappa^{!}=0\) on this character says \(\chi_{s}\) and
\(\chi_{1-s}\) are cohomologically identified, equivalent to the
functional-equation identity~\eqref{v4-arith:w5-direct:eq:functional-equation}.
\end{remark}

\subsection{Archimedean Heisenberg Character}
\label{v4-arith:w5-direct:step-2}

\begin{definition}[archimedean Heisenberg character]
\label{v4-arith:w5-direct:def:zhu}
The archimedean fibre \(\mathsf{H}_{\infty}\) is the standard Heisenberg
vertex operator algebra on \(\mathbb{C}_{\infty}\) with Virasoro zero-mode
\(L_{0}\) and central charge \(c_{\infty}=1\) (one free boson). For the
primary \(\mathsf{H}_{\infty}\)-module \(V_{\infty}\) (the vacuum
representation), the Zhu character is
\begin{equation}
\label{v4-arith:w5-direct:eq:zhu-char}
\chi_{V_{\infty}}(\tau)\;:=\;\mathrm{tr}_{V_{\infty}}\bigl(q^{L_{0}-c_{\infty}/24}\bigr),
\qquad q=e^{2\pi i\tau},
\end{equation}
a standard VOA Zhu trace on the archimedean fibre. The global object
\(\mathcal{V}^{\mathrm{prim}}\) is the Hochschild trace class of the
identity on the arithmetic Hecke category; it is \emph{not} a vertex
algebra. Identifying its Hochschild character with \(\Lambda_{\zeta}(s)\)
requires a native trace comparison with the adelic zeta integral,
as specified in \Cref{v4-arith:prim-zeta:native-character-hypothesis}.
\end{definition}

\begin{theorem}[P2: Hochschild character and functional equation]
\label{v4-arith:w5-direct:thm:P2-elementary}
\ClaimStatusProvedHereConditional. Assume the native Hochschild-to-zeta
trace comparison, with its actual coefficient module and domains. The Hochschild character of
\(\mathcal{V}^{\mathrm{prim}}\) equals \(\Lambda_{\zeta}(s)\) under the
identification of the Hochschild trace with Tate's adelic zeta integral
(\Cref{v4-arith:thm:G-row-Riemann-functional-equation}). The primary
determinant carries rank residue \(r_{\zeta}=2\), the residue at
\(s=1\) of \(2\zeta(s)\) (\Cref{v4-arith:thm:c-as-residue}). The adelic
functional equation \(\Lambda_{\zeta}(s)=\Lambda_{\zeta}(1-s)\) follows
from Poisson summation applied to the Schwartz--Bruhat function at each
place.  The Heisenberg vacuum character is the local eta-factor; the
theta transformation enters only after the lattice Schwartz sector is
adjoined.
\end{theorem}

\begin{remark}[elementary content of P2]
\label{v4-arith:w5-direct:rem:P2-elementary}
On the archimedean fibre,
\[
\chi_{V_{\infty}}(\tau)=q^{-1/24}\prod_{n\geq1}(1-q^{n})^{-1}.
\]
This eta-character is not Tate's theta integral.  Tate's local factor is
obtained from the Schwartz function on \(\mathbb{R}\) and the theta sum
\(\theta(t)=\sum_{n\in\mathbb{Z}}e^{-\pi n^{2}t}\); Poisson summation is
the identity \(\theta(t)=t^{-1/2}\theta(1/t)\).  No finite Zhu
\(S\)-matrix or rational-VOA modularity is used in the Hochschild trace.
Tate (1950) identified \(\Lambda_{\zeta}\) with the adelic zeta integral;
a native trace-preserving map is required to transport that scalar computation to a Hochschild-trace
identity on \(\mathcal{V}^{\mathrm{prim}}\), with VOA language confined
to \(\mathsf{H}_{\infty}\). Source: J.~T.~Tate, thesis, Princeton 1950
(reprinted in Cassels--Fr\"{o}hlich, ``Algebraic Number Theory'',
Academic Press 1967).
\end{remark}

\subsection{Step 3: Petersson positivity}
\label{v4-arith:w5-direct:step-3}

\begin{definition}[adelic Petersson form]
\label{v4-arith:w5-direct:def:petersson}
For \(\phi_{1},\phi_{2}\in\mathcal{V}^{\mathrm{prim}}\) with finite
adelic support, define
\begin{equation}
\label{v4-arith:w5-direct:eq:petersson}
\langle \phi_{1},\phi_{2}\rangle_{\mathrm{Pet}}\;:=\;
\int_{[GL_{2}]}\phi_{1}(h)\,\overline{\phi_{2}(h)}\,dh,
\end{equation}
the \(L^{2}\)-inner product on the adelic quotient
\([GL_{2}]:=GL_{2}(\mathbb{Q})\backslash GL_{2}(\mathbb{A})\) with
respect to the Tamagawa measure \(dh\). Equivalently, the restricted tensor
decomposition of \(\mathcal{V}^{\mathrm{prim}}\) yields
\begin{equation}
\label{v4-arith:w5-direct:eq:petersson-factorized}
\langle \phi_{1},\phi_{2}\rangle_{\mathrm{Pet}}\;=\;
\langle (\phi_{1})_{\infty},(\phi_{2})_{\infty}\rangle_{\infty}
\;\cdot\;{\prod_{p}}'\langle (\phi_{1})_{p},(\phi_{2})_{p}\rangle_{p},
\end{equation}
the restricted product over places of local Petersson forms.
\end{definition}

\begin{theorem}[Petersson positivity and its native pullback]
\label{v4-arith:w5-direct:thm:P3-elementary}
\ClaimStatusProvedHereConditional.
Suppose the stated native cuspidal core admits an injective linear map
into the cuspidal \(L^2\) space, and the local and adelic pairing maps
of \Cref{v4-arith:kp:hyp:native-pairing-maps} exist. The pulled-back
Petersson form is positive definite and equals the native bilinear
Koszul form after the conjugate-linear duality. These map hypotheses
apply also to the holomorphic and residual sectors. Temperedness or
Ramanujan bounds alone do not construct the maps.
\end{theorem}
\begin{proof}
For an injective map \(I\), the integral \(\|Iv\|_{L^2}^2\) is positive
for \(v\ne0\). The identity of forms follows from
\Cref{v4-arith:thermal:pairing-transfer}; its local hypotheses include
both forms and the conjugate-linear map between their second carriers.
\end{proof}

\begin{remark}[the concrete radial sector]
\label{v4-arith:w5-direct:rem:P3-elementary}
In \Cref{v4-arith:kp:radial-pairing-realization} the shell functions
\(p^{-r}\mathbf1_{p^r\mathbb Z_p^\times}\), with measure
\(|a|_p^{-1}d^\times a\) and unit-group volume one, have squared norm
\(p^{-r}\). Their maps to creation and division vectors prove the
pairing identity on this radial sector. Extending those maps to the
native automorphic and Koszul carriers, with their phases and Hecke
operations, is an additional problem.
\end{remark}

\subsection{Step 4: Bost--Connes and the Hilbert--P\'olya Boundary}
\label{v4-arith:w5-direct:step-4}

\begin{definition}[the represented prime thermal system]
\label{v4-arith:w5-direct:def:BC}
For each prime \(p\), let \(S_p\) be the unilateral isometry on
\(\ell^2(\mathbb N_0)\), \(R_p=S_p^*\), and let
\(\alpha_t(S_p)=p^{it}S_p\). For real \(\beta>0\), the local state is
\[
 \Phi_{p,\beta}(a)=(1-p^{-\beta})\sum_{r\ge0}p^{-\beta r}
                 \langle v_r,av_r\rangle.
\]
The product states on finite tensor products agree on the algebraic
union. Its completion in the GNS norm at \(\beta=1\) is
\(\mathcal H_1\) of \Cref{v4-arith:w5-b:def:gns}.
These thermal states satisfy the KMS identity, which is not traciality:
\(\Phi(R_pS_p)=1\) but \(\Phi(S_pR_p)=p^{-\beta}\).
An arithmetic algebra containing all phase operators requires its
own compatible representation. The occupation representation on
\(\ell^2(\mathbb N^\times)\) has Gibbs partition function
\(\zeta(\beta)\) only for \(\beta>1\).
\end{definition}

\begin{theorem}[modular and occupation normalization]
\label{v4-arith:w5-direct:thm:HPAr-elementary}
\ClaimStatusProvedHere.
On the product GNS space at \(\beta=1\),
\(N_{\mathrm{BC}}=-\tfrac12\log\Delta_1\) is self-adjoint. Its
point eigenvalues are \(\tfrac12\log(a/b)\), \(a,b\in\mathbb N^\times\),
each with infinite multiplicity; its spectrum is \(\mathbb R\).
On the closed creation subspace there is a unitary
\(W:\ell^2(\mathbb N^\times)\to\mathcal H_1^+\),
\(Wv_m=[\mu_m]\), such that
\[
        WN_{\mathrm{occ}}=2N_{\mathrm{BC}}W,
        \qquad N_{\mathrm{occ}}v_m=(\log m)v_m.
\]
\end{theorem}
\begin{proof}
The local modular eigenvalues on matrix units are \(p^{j-i}\).
Their finite tensor products give the stated eigenvalues and a complete
orthonormal eigenbasis. The creation vectors have norm one and modular
eigenvalue \(m^{-1}\). The maximal diagonal operator domains give the
intertwining identity; see \Cref{v4-arith:w5-b:occupation-comparison}.
\end{proof}

\begin{remark}[the zero-divisor problem]
\label{v4-arith:w5-direct:rem:HPAr-stop}
These operators construct occupation weights. A self-adjoint operator
whose determinant has divisor \(\mathrm{div}\,\xi_{\mathbb R}\)
requires further arithmetic and trace maps. Literal injective
intertwining of \(N_{\mathrm{occ}}\) with that operator on the vacuum
is impossible: it would force a zero at \(s=1/2\), whereas
\(\xi_{\mathbb R}(1/2)\ne0\)
(\Cref{v4-arith:w5-b:vacuum-obstruction}). The grading and the proposed
zero-divisor operator in \Cref{v4-arith:w5-b:def:twist} are separate data.
\end{remark}

\subsection{Step 5: Weil positivity via Mellin--Petersson isometry}
\label{v4-arith:w5-direct:step-5}

\begin{definition}[Mellin lift]
\label{v4-arith:w5-direct:def:mellin-lift}
For \(f\in\mathrm{PW}(\mathbb{R})\) in a suitable test class (compactly
supported and even), define
\begin{equation}
\label{v4-arith:w5-direct:eq:mellin-lift}
\Upsilon(f)(z)\;:=\;\sum_{n=1}^{\infty}f(\log n)\,n^{1/2}\,e^{2\pi inz},
\qquad z\in\mathbb{H}.
\end{equation}
The series converges on the upper half plane to a holomorphic function
of rapid decay at the cusp. A native cuspidal or automorphic interpretation of this holomorphic
series requires an additional realization map. Its Fourier expansion
alone does not prove modular transformation laws or native membership.
\end{definition}

\begin{hypothesis}[a native Weil--Petersson norm realization]
\label{v4-arith:w5-direct:thm:WP-isometry}
On a specified arithmetic test space \(\mathcal E\), assume a map
\(\Upsilon:\mathcal E\to\mathcal V^{\mathrm{prim}}_{\mathrm{cusp}}\)
whose image lies in an injectively realized positive Petersson core,
and assume the actual norm identity
\begin{equation}
 W(f*\widetilde f)=\langle\Upsilon(f),\Upsilon(f)\rangle_{\mathrm{Pet}}.
 \label{v4-arith:w5-direct:eq:WP-isometry}
\end{equation}
The Fourier series above is a proposed construction to test against this
identity; native membership and the identity are separate requirements.
\end{hypothesis}

\begin{theorem}[positivity under a realized norm identity]
\label{v4-arith:w5-direct:thm:HPLi-elementary}
Under the norm realization hypothesis, \(W(f*\widetilde f)\ge0\) for
\(f\in\mathcal E\). If a specified Li trace is a finite limit of these
nonnegative values along test functions in \(\mathcal E\), that trace
is also nonnegative. The approximation and trace identity are additional
hypotheses; this statement does not assert them for all Li coefficients.
\end{theorem}
\begin{proof}
The squared Petersson norm is nonnegative on its realized Hilbert core.
Substitution in the assumed identity proves the first assertion. A
finite limit of nonnegative real numbers is nonnegative, proving the
second. Neither assertion constructs the native map or the Li trace
approximation.
\end{proof}

\section{The Restricted Test Class}
\label{v4-arith:w5-direct:closing-restricted}

\begin{theorem}[restricted-class Weil positivity]
\label{v4-arith:w5-direct:thm:restricted-RH}
\ClaimStatusProvedHere. Assume the domain-restricted P3
(\Cref{v4-arith:w5-direct:thm:P3-elementary}), the Weil--Petersson
isometry on the restricted test class
(\Cref{v4-arith:w5-direct:thm:WP-isometry}), and HP--Li positivity on
the restricted test class
(\Cref{v4-arith:w5-direct:thm:HPLi-elementary}). Then
\[
W(f\ast\widetilde f)\geq0
\]
for every \(f\) in the restricted test class.  No zero-placement
conclusion follows until the positivity is extended to the full
admissible Weil class.
\end{theorem}

\begin{proof}
The domain-restricted P3 and the Weil--Petersson isometry give
\[
W(f\ast\widetilde f)
=\langle\Upsilon(f),\Upsilon(f)\rangle_{\mathrm{Pet}}
\]
on the restricted domain.  The right hand side is non-negative by
P3.  Weil's RH criterion applies only after this inequality holds for
all admissible Weil test functions, so the restricted statement is not
a zero-placement theorem.
\end{proof}

\begin{remark}[domain boundary]
\label{v4-arith:w5-direct:rem:what-not-proved}
\Cref{v4-arith:w5-direct:thm:restricted-RH} is a theorem on a
proper test domain.  The passage to the full Weil cone requires:
\begin{itemize}
\item[(i)] P3 on the non-tempered and Eisenstein sectors
(\Cref{v4-arith:w5-c:thm:main});
\item[(ii)] extending the Weil--Petersson isometry to an unrestricted
test class compatible with all admissible Weil test functions
(\Cref{v4-arith:w5-a:sec:obstruction});
\item[(iii)] handling the Eisenstein continuous-spectrum contribution
either by domain restriction or Langlands--Shahidi normalization
correction
(\Cref{v4-arith:w5-c:eisenstein}).
\end{itemize}
These three analytic assertions are independent of positivity on the
restricted class.
\end{remark}

\section{Restricted-Class Comparison}
\label{v4-arith:w5-direct:construction~5-map}

The restricted-class comparison partitions the conditions of
\Cref{v4-arith:w5-direct:rem:what-not-proved} into disjoint
mathematical reductions.

\begin{center}
\small
\begin{tabular}{p{0.45\textwidth}|p{0.50\textwidth}}
Condition & Reduction \\
\hline
F1.a.1 Banerjee--Clozel ideal &
\Cref{v4-arith:w5-e:chapter} \\
F1.a.2 small-prime extension (\(p=2,3\)) &
\Cref{v4-arith:w5-e:chapter} \\
F1.a.3.A block-Hochschild comparison &
\Cref{v4-arith:w5-f:F1a3-AB} \\
F1.a.3.B Paskunas \(\mathrm{Ext}^{1}\) vanishing &
\Cref{v4-arith:w5-f:F1a3-AB} \\
F1.a.3.C wild-EG compact generation (\(p=2,3\)) &
\Cref{v4-arith:w5-g:f1a3-C} \\
\(\mathrm{L}_{\mathrm{ACD}}\) adelic categorical descent &
\Cref{v4-arith:w5-e:chapter} \\
BC-diamond-glue (Ch1) Chenevier pseudochar &
\Cref{v4-arith:w5-h:bcglue-core-and-F2c} \\
BC-diamond-glue (AG.det) Arinkin--Gaitsgory split &
\Cref{v4-arith:w5-h:bcglue-core-and-F2c} \\
F2.c restricted-tensor \(E_{2}\) transport &
\Cref{v4-arith:w5-h:bcglue-core-and-F2c} \\
Arith.~Positselski: continuous duals &
\Cref{v4-arith:w5-d:chapter} \\
Arith.~Positselski: derived completion &
\Cref{v4-arith:w5-d:chapter} \\
Arith.~Positselski: faithful lifting &
\Cref{v4-arith:w5-d:chapter} \\
Arith.~Positselski: Koszul-dual BC &
\Cref{v4-arith:w5-d:chapter} \\
Full \(\mathcal{V}^{\mathrm{prim}}\) P3: non-tempered (Ramanujan) &
\Cref{v4-arith:w5-c:full-vprim-P3} \\
Full \(\mathcal{V}^{\mathrm{prim}}\) P3: Eisenstein density correction &
\Cref{v4-arith:w5-c:full-vprim-P3} \\
H-P\(^{\mathrm{Ar}}\) (W) weight existence &
\Cref{v4-arith:w5-b:hp-operator} \\
(H) Weil formula at Gaussians &
\Cref{v4-arith:w5-i:chapter} \\
R\(_{\mathrm{analytic}}\) density-of-zeros frontier &
\Cref{v4-arith:w5-i:chapter} \\
R1.formal.b/c formal-Noetherian residuals &
\Cref{v4-arith:w5-i:chapter} \\
VB\(^{\ast}\) main direction primary Fock &
\Cref{v4-arith:w5-j:VBstar-spectral-positivity} \\
HP--Li restricted-class positivity &
\Cref{v4-arith:w5-a:chapter} \\
\end{tabular}
\end{center}

\section{The Full Weil Cone}
\label{v4-arith:w5-direct:irreducible-residual}

The preceding reductions leave a single full-domain statement.

\begin{theorem}[full-domain Weil positivity]
\label{v4-arith:w5-direct:thm:irreducible-residual}
\ClaimStatusProvedHere. Assume the restricted-class comparison
conditions of \Cref{v4-arith:w5-direct:construction~5-map} hold. Then RH is equivalent to
the unrestricted Weil-positivity statement:
\begin{equation}
\label{v4-arith:w5-direct:eq:irreducible}
W(f\ast\widetilde f)\;\geq\;0\quad\text{for every admissible Weil test function }f.
\end{equation}
Equivalently, RH is equivalent to the full-domain extension of
\Cref{v4-arith:w5-direct:thm:HPLi-elementary} from the restricted test
class to the unrestricted Paley--Wiener class.
\end{theorem}

\begin{proof}
The assumptions give the restricted comparisons on their stated
domains.  Weil's theorem,
\Cref{v4-arith:w5-direct:thm:weil-positivity}, says that positivity of
the Weil distribution on the full admissible Paley--Wiener class is
equivalent to RH.  Thus the full-domain extension is neither weaker
nor stronger than RH; it is RH in Weil's form.
\end{proof}

\begin{remark}[domain of the reduction]
\label{v4-arith:w5-direct:rem:reduction-honesty}
\Cref{v4-arith:w5-direct:thm:irreducible-residual} identifies the
unrestricted extension with RH.  The restricted construction supplies
no theorem forcing this extension; the continuous and Eisenstein
sectors remain part of the full Weil inequality.
\end{remark}

\section{Classical Boundary}
\label{v4-arith:w5-direct:recap}

\begin{enumerate}
\item \textbf{RH classical form}
(\Cref{v4-arith:w5-direct:conj:RH}): zeros of \(\xi_{\mathbb{R}}\) on
\(\mathrm{Re}(s)=1/2\).
\item \textbf{Four classical equivalents}
(\Cref{v4-arith:w5-direct:rem:four-equivalents}): Weil positivity,
Li positivity, Hilbert--P\'olya, de Branges.
\item \textbf{Five structural ingredients}:
Heisenberg + Koszul (P1), Hochschild character + functional equation
(P2), Petersson positivity (P3), Bost--Connes log-modular operator,
Mellin--Petersson lift.
\item \textbf{Restricted-class positivity}
(\Cref{v4-arith:w5-direct:thm:restricted-RH}): Weil positivity on the
restricted test class.
\item \textbf{Restricted-class comparison}
(\Cref{v4-arith:w5-direct:construction~5-map}): the comparison
conditions factor through ten named mathematical reductions.
\item \textbf{Full Weil cone}
(\Cref{v4-arith:w5-direct:thm:irreducible-residual}): extension of
the Weil--Petersson isometry to the unrestricted Paley--Wiener
class, equivalent to RH.
\end{enumerate}

Each step in the construction above is independent of Vols~I--III
and has a stated mathematical domain. The full-domain statement
(\Cref{v4-arith:w5-direct:thm:irreducible-residual}) is the Weil--Petersson
isometry extension to the unrestricted Paley--Wiener class.
